\section{全同粒子}

Gross 第 1 章强调：经典可区分的全同粒子在量子力学中不可区分，波函数必须按交换对称性分类。

\subsection{交换对称性}

自然界中的基本粒子（及许多准粒子）是全同的：交换任意两个粒子，物理上不可区分。量子力学要求多体波函数在交换下要么对称（玻色子），要么反对称（费米子）：
\begin{equation}
  P_{ij}\Psi=\pm\Psi.
\end{equation}
费米子的反对称性导致 Pauli 不相容原理：两个费米子不能占据同一单粒子态。自旋--统计定理把自旋与对称性绑定（半整数自旋 $\Rightarrow$ 费米子）。

\subsection{Slater 行列式与永久式}

$N$ 个费米子若占据正交单粒子轨道 $\{\phi_{k_1},\ldots,\phi_{k_N}\}$，则多体波函数可取 Slater 行列式
\begin{equation}
  \Psi(\mathbf{x}_1,\ldots,\mathbf{x}_N)
  =\frac{1}{\sqrt{N!}}
  \det\bigl[\phi_{k_i}(\mathbf{x}_j)\bigr],
\end{equation}
其中 $\mathbf{x}=(\mathbf{r},\sigma)$。玻色子对应的是永久式（permanent）。Hartree--Fock 理论正是在 Slater 行列式流形上做变分优化。

\subsection{约化密度矩阵（预告）}

对纯态 $|\Psi\rangle$，一体约化密度矩阵
\begin{equation}
  \rho_1(\mathbf{x},\mathbf{x}')
  =N\int\mathrm{d}\mathbf{x}_2\cdots\mathrm{d}\mathbf{x}_N
  \Psi(\mathbf{x},\mathbf{x}_2,\ldots)
  \Psi^*(\mathbf{x}',\mathbf{x}_2,\ldots)
\end{equation}
的本征值是自然轨道占据数。二次量子化后 $\rho_1=\langle\psi^\dagger\psi\rangle$，见密度算符一节。

\subsection{粒子数表象}

与其跟踪坐标空间中的巨大波函数，不如用各单粒子态的占据数 $\{n_\alpha\}$ 标记态。对费米子 $n_\alpha=0,1$，对玻色子 $n_\alpha=0,1,2,\ldots$。全体占据数组态张成 Fock 空间
\begin{equation}
  \mathcal{F}
  =\bigoplus_{N=0}^\infty \mathcal{H}_N^{(\pm)},
\end{equation}
其中 $\mathcal{H}_N^{(\pm)}$ 为 $N$ 粒子对称/反对称子空间。二次量子化算符在 $\mathcal{F}$ 上自然作用。

\begin{note}
真空态 $|0\rangle$ 表示所有单粒子态均未被占据。多体激发可理解为在真空上作用一串产生算符得到的态。对有限密度体系，常以填满的费米海为“真空”，其上的粒子与空穴激发构成低能自由度。
\end{note}
